% ============================================================================
%  SECCIÓN 4.8: FORMULACIÓN Y VALIDACIONES
% ============================================================================

\section{Formulaci\'on y validaciones}

Las validaciones garantizan que los datos ingresados por el usuario cumplan
con los requisitos del sistema antes de ser procesadas. En este proyecto,
las validaciones se implementan mediante esquemas Zod que definen reglas
claras para cada campo.

\subsection{Reglas de validaci\'on}

Las reglas de validaci\'on se establecen de acuerdo con los requerimientos
funcionales del sistema:

\begin{table}[H]
  \centering
  \caption{Reglas de validaci\'on de formularios}
  \contenidoConFuente{%
    \begin{tablaAPA}
      \begin{tabular}{@{}p{3.4cm}p{5.6cm}p{6.6cm}@{}}
        \toprule
        \textbf{Campo} & \textbf{Regla} & \textbf{Mensaje de error} \\
        \midrule
        Nombre completo & Obligatorio, m\'inimo 3 caracteres &
        ``El nombre debe tener al menos 3 caracteres''. \\
        \addlinespace
        Email & Obligatorio, formato v\'alido &
        ``Ingresa un email v\'alido''. \\
        \addlinespace
        Contrase\~na & Obligatorio, m\'inimo 6 caracteres &
        ``La contrase\~na debe tener al menos 6 caracteres''. \\
        \addlinespace
        Confirmar contrase\~na & Debe coincidir con contrase\~na &
        ``Las contrase\~nas no coinciden''. \\
        \bottomrule
      \end{tabular}
    \end{tablaAPA}%
  }{Elaboraci\'on propia en base a los esquemas Zod del proyecto.}
\end{table}

\subsection{Implementaci\'on de validaciones con Zod}

Cada formulario define un esquema Zod que codifica las reglas de
validaci\'on. El esquema del formulario de registro valida cuatro campos
y verifica la coincidencia de contrase\~nas mediante \textit{refine}:

\begin{lstlisting}
const registerSchema = z
  .object({
    fullName: z
      .string()
      .min(1, 'El nombre es obligatorio')
      .min(3, 'El nombre debe tener al menos 3 caracteres')
      .max(100, 'El nombre es demasiado largo'),
    email: z
      .string()
      .min(1, 'El email es obligatorio')
      .email('Ingresa un email valido'),
    password: z
      .string()
      .min(1, 'La contrasena es obligatoria')
      .min(6, 'La contrasena debe tener al menos 6 caracteres'),
    confirmPassword: z
      .string()
      .min(1, 'Confirma tu contrasena'),
  })
  .refine(
    (data) => data.password === data.confirmPassword,
    {
      message: 'Las contrasenas no coinciden',
      path: ['confirmPassword'],
    }
  );
\end{lstlisting}

\subsection{Visualizaci\'on de errores}

Los errores de validaci\'on se muestran debajo de cada campo en tiempo real,
mientras el usuario escribe. El borde del campo cambia a color rojo para
se\~nalar visualmente el error:

\begin{lstlisting}
<TextInput
  style={[styles.input, errors.email && styles.inputError]}
  placeholder="tu@email.com"
  value={value}
  onChangeText={onChange}
  onBlur={onBlur}
/>
{errors.email && (
  <Text style={styles.errorText}>{errors.email.message}</Text>
)}
\end{lstlisting}

Los estilos aplicados son:

\begin{lstlisting}
inputError: {
  borderColor: colors.error,  // Rojo #EF4444
},
errorText: {
  fontSize: fontSize.xs,
  color: colors.error,
  marginTop: spacing.xs,
},
\end{lstlisting}

\subsection{Errores del servidor}

Adem\'as de las validaciones locales, la aplicaci\'on maneja errores
provenientes del servidor (Supabase), como credenciales inv\'alidas o
correos ya registrados. Estos errores se muestran en una caja especial
con fondo rojo claro:

\begin{lstlisting}
{serverError && (
  <View style={styles.serverErrorBox}>
    <Text style={styles.serverErrorText}>{serverError}</Text>
  </View>
)}
\end{lstlisting}

\subsection{Control del estado de los botones}

Los botones de env\'io se deshabilitan durante el proceso de
autenticaci\'on para evitar env\'ios duplicados. Mientras se procesa la
solicitud, se muestra un indicador de carga (\textit{ActivityIndicator})
en lugar del texto del bot\'on:

\begin{lstlisting}
<TouchableOpacity
  style={[styles.button, isSubmitting && styles.buttonDisabled]}
  onPress={handleSubmit(onSubmit)}
  disabled={isSubmitting}
>
  {isSubmitting ? (
    <ActivityIndicator color={colors.white} />
  ) : (
    <Text style={styles.buttonText}>Iniciar sesion</Text>
  )}
</TouchableOpacity>
\end{lstlisting}

El estilo \textit{buttonDisabled} reduce la opacidad del bot\'on al 60\\%
para indicar visualmente que no est\'a disponible.
